\documentclass[11pt]{article}
% ============================================================================
% preamble.tex — Estilo común de los manuales de práctica SISELA (FIME / UANL)
%
% Reproduce (de forma aproximada) el estilo de los PDF originales
% (~/Descargas/manuales/Manuales/pN.pdf), de los que no existe fuente .tex.
% Compilar con:  pdflatex  (o  latexmk -pdf pN.tex)
% ============================================================================

\usepackage[utf8]{inputenc}
\usepackage[T1]{fontenc}
\usepackage[spanish,es-tabla]{babel}
\usepackage[a4paper,margin=2.5cm,headheight=15pt]{geometry}

% --- Tipografía: serif para el cuerpo, sans para los títulos ---
\usepackage{libertinus}
\usepackage[scaled=0.92]{helvet}   % sans para encabezados
\usepackage{microtype}

% --- Matemáticas ---
\usepackage{amsmath,amssymb}
\usepackage{siunitx}
\sisetup{detect-all, per-mode=symbol}

% --- Tablas y figuras ---
\usepackage{booktabs}
\usepackage{array}
\usepackage{multirow}
\usepackage{caption}
\captionsetup{font=small,labelfont=bf,labelsep=period}
\usepackage{graphicx}
\usepackage{float}

% --- Diagramas ---
\usepackage{tikz}
\usetikzlibrary{shapes.geometric,arrows.meta,positioning,calc}
\tikzset{
  block/.style   = {draw, rectangle, minimum width=2.6cm, minimum height=0.9cm,
                    align=center, font=\footnotesize},
  term/.style    = {draw, rounded corners, minimum width=2.2cm, minimum height=0.8cm,
                    align=center, font=\footnotesize, fill=black!5},
  decision/.style= {draw, diamond, aspect=2, align=center, font=\scriptsize,
                    inner sep=1pt},
  flow/.style    = {-{Latex[length=2mm]}, thick},
}

% --- Código ---
\usepackage{xcolor}
\definecolor{codebg}{gray}{0.96}
\definecolor{codekw}{rgb}{0.13,0.29,0.53}
\definecolor{codecm}{rgb}{0.40,0.45,0.40}
\usepackage{listings}
\lstdefinestyle{sisela}{
  backgroundcolor=\color{codebg}, basicstyle=\ttfamily\footnotesize,
  keywordstyle=\color{codekw}\bfseries, commentstyle=\color{codecm}\itshape,
  stringstyle=\color{codekw}, numbers=left, numberstyle=\tiny\color{gray},
  numbersep=6pt, frame=none, breaklines=true, showstringspaces=false,
  columns=fullflexible, xleftmargin=1.8em, aboveskip=1em, belowskip=1em,
  literate={á}{{\'a}}1 {é}{{\'e}}1 {í}{{\'i}}1 {ó}{{\'o}}1 {ú}{{\'u}}1
           {ñ}{{\~n}}1 {Á}{{\'A}}1 {°}{{\textdegree}}1 {µ}{{$\mu$}}1
           {✓}{{\checkmark}}1,
}
\lstset{style=sisela}
\newcommand{\code}[1]{\texttt{\small #1}}

% --- Listas ---
\usepackage{enumitem}
\setlist{nosep, leftmargin=1.6em}

% --- Encabezado / pie de página ---
\usepackage{fancyhdr}
\usepackage{lastpage}
\pagestyle{fancy}
\fancyhf{}
\fancyhead[L]{\footnotesize\sffamily Sistemas Eléctricos de Aeronaves — Práctica \practnum}
\fancyhead[R]{\footnotesize\sffamily FIME / UANL}
\fancyfoot[C]{\footnotesize P\'agina \thepage\ de \pageref{LastPage}}
\renewcommand{\headrulewidth}{0.4pt}
\renewcommand{\footrulewidth}{0.4pt}

% --- Títulos de sección en sans-serif ---
\usepackage{titlesec}
\titleformat{\section}{\normalfont\Large\sffamily\bfseries}{\thesection}{0.7em}{}
\titleformat{\subsection}{\normalfont\large\sffamily\bfseries}{\thesubsection}{0.6em}{}
\titleformat{\subsubsection}{\normalfont\normalsize\sffamily\bfseries}{\thesubsubsection}{0.5em}{}

% --- hyperref al final ---
\usepackage[hidelinks,pdfusetitle]{hyperref}

% ============================================================================
% Comandos de portada y cajas
% ============================================================================
\newcommand{\practnum}{N}          % lo redefine cada pN.tex
\newcommand{\practtitle}{}         % idem

\newcommand{\portada}{%
  \thispagestyle{empty}
  \begin{center}
    {\large Universidad Autónoma de Nuevo León}\\[2pt]
    {\normalsize Facultad de Ingeniería Mecánica y Eléctrica}\\[2.2cm]
    \rule{\linewidth}{0.4pt}\\[0.4cm]
    {\Huge\sffamily\bfseries Práctica \practnum}\\[0.5cm]
    {\large \practtitle}\\[0.3cm]
    \rule{\linewidth}{0.4pt}\\[2.5cm]
    {\normalsize Laboratorio de Sistemas Electrónicos de Aeronaves\\(SISELA)}\\[1.5cm]
    {\normalsize Ciclo Escolar 2026}\\[3cm]
  \end{center}
}

\newenvironment{resumen}{\par\noindent\rule{\linewidth}{0.4pt}\par\smallskip
  \small\itshape\noindent\textbf{Resumen Ejecutivo} ---\ }{\par\smallskip
  \noindent\rule{\linewidth}{0.4pt}\par}

\newenvironment{nota}[1][Nota]{\par\smallskip\noindent\rule{\linewidth}{0.4pt}\par\smallskip
  \small\itshape\noindent\textbf{#1} ---\ }{\par\smallskip
  \noindent\rule{\linewidth}{0.4pt}\par\smallskip}

\renewcommand{\practnum}{8}
\renewcommand{\practtitle}{Integración de Sistemas: Adquisición, ARINC 429 y Actuación\\
Análisis de Señales y Módulo Opcional de Dron — Énfasis Aeronáutico}

\begin{document}
\portada

\begin{resumen}
Los sistemas aeronáuticos modernos operan bajo una premisa fundamental: múltiples
subsistemas independientes deben intercambiar datos de manera confiable, determinista y
verificable. Esta práctica integra los conocimientos de las Prácticas 1--7 para construir
un sistema completo de adquisición, codificación ARINC 429 y actuación. El manual
conserva las \textbf{cuatro variantes} originales (velocidad, temperatura, altitud,
superficie de vuelo FBW) como marco teórico de diseño, y documenta la
\textbf{implementación de referencia del repositorio}: un panel integrado de
instrumentos (sensores + servos + motor + tren de aterrizaje) con nueve modos, que
incorpora además \textbf{análisis de señales} (latencia de la cadena de control,
muestreo multicanal) y un \textbf{módulo opcional} para construir un cuadricóptero con
motores brushless 2212.
\end{resumen}

\tableofcontents
\newpage

% ============================================================================
\section{Introducción}

Los sistemas aeronáuticos modernos operan bajo una premisa fundamental: múltiples
subsistemas independientes deben intercambiar datos de manera confiable, determinista y
verificable. Desde la lectura de un sensor de temperatura en la bahía de equipos hasta el
posicionamiento de una aguja indicadora en cabina, cada dato atraviesa un proceso riguroso
de adquisición, codificación, transmisión, validación y actuación.

El estándar ARINC 429 (\emph{Mark 33 Digital Information Transfer System}) define el
protocolo de comunicación de datos más extendido en la aviación comercial. Cada palabra
de 32 bits transporta una variable identificada por su \emph{Label}, con campos de estado
(SSM), datos escalados (\emph{Data Field}) y un bit de paridad impar que garantiza la
integridad de la trama. El bus opera de forma unidireccional: un transmisor (\emph{source})
envía a uno o más receptores (\emph{sinks}), lo que simplifica la arquitectura y elimina
colisiones.

Esta práctica integra los conocimientos de las Prácticas 1--7 para construir un sistema
completo de adquisición y presentación de datos aeronáuticos. El planteamiento original
propone que cada equipo seleccione una de \textbf{cuatro variantes} y desarrolle dos
fases: (1) \textbf{Fase de Adquisición} (Nodo Sensor) — un MCU lee un sensor físico,
convierte la señal eléctrica a un valor físico y lo codifica en una palabra ARINC 429 de
32 bits; (2) \textbf{Fase de Salida} (Nodo Indicador) — un segundo MCU (o el mismo MCU en
configuración de lazo) recibe la palabra ARINC 429, la decodifica, valida su integridad y
comanda un actuador para presentar la información.

\begin{nota}[Implementación de referencia del repositorio]
El repositorio SISELA-Init implementa P8 como un \textbf{panel integrado} (un solo
programa con 9 modos: instrumentos, control manual, potencia, tren, piloto automático,
diagnóstico, configuración, dron opcional y análisis de señales) en lugar de 4 proyectos
de variante independientes. Los conceptos ARINC 429 de esta sección (BNR, SSM, paridad)
son directamente transferibles: cada equipo puede optar por (a) seguir el panel integrado
del repositorio, añadiendo la codificación/decodificación ARINC 429 sobre sus funciones
de \code{lib/sensors.py} y \code{lib/flight\_controls.py} existentes, o (b) implementar
una de las 4 variantes originales como proyecto independiente siguiendo la \S7 de este
manual. Ambos caminos cubren los mismos objetivos de aprendizaje.
\end{nota}

% ============================================================================
\section{Objetivos de Aprendizaje}
\begin{enumerate}
  \item Diseñar un sistema integrado que combine adquisición de datos, codificación
        digital y actuación coordinada.
  \item Implementar la estructura de la palabra ARINC 429 de 32 bits (Label, SDI, Data,
        SSM, Parity) en software embebido.
  \item Codificar valores físicos en formato BNR (\emph{Binary Number Representation})
        con escalado configurable.
  \item Calcular y verificar el bit de paridad impar para garantizar la integridad de la
        trama.
  \item Utilizar la Matriz de Estado del Signo (SSM) para reportar el estado del dato.
  \item Transmitir y recibir palabras ARINC 429 entre nodos usando UART como medio físico
        simplificado.
  \item Decodificar la palabra recibida y convertir el dato BNR de vuelta a unidades
        físicas.
  \item Comandar actuadores (\emph{stepper}, servo, LEDs, \emph{buzzer}) de forma
        proporcional al dato decodificado.
  \item \textbf{Medir la latencia de la cadena sensor$\to$actuador} y el comportamiento
        del sistema al muestrear varios canales analógicos simultáneamente.
  \item \textbf{(Opcional) Integrar motores brushless 2212} mediante ESC, con una
        máquina de armado \emph{fail-safe} y un mezclador de cuadricóptero en X, como
        base de un proyecto de extensión hacia un dron real.
  \item Documentar el sistema completo bajo el formato de reporte AIAA.
  \item Relacionar la práctica con sistemas reales de aviónica: FMS, EICAS, EFIS y sus
        interfaces de bus.
\end{enumerate}

% ============================================================================
\section{Fundamentos Teóricos}

\subsection{Protocolo ARINC 429}
ARINC 429 es un bus de datos serial unidireccional utilizado en la mayoría de las
aeronaves comerciales (Boeing 737/747/757/767/777, Airbus A300/A310/A320, entre otras).
Sus características principales son: bus unidireccional con un único transmisor y hasta
20 receptores, velocidad de 12.5\,kbit/s (baja) o 100\,kbit/s (alta), medio físico de par
trenzado blindado con señalización diferencial ($\pm\SI{10}{\volt}$ en sistemas reales),
y palabras de 32 bits transmitidas LSB primero.

\subsection{Estructura de la palabra ARINC 429 (32 bits)}
\begin{table}[H]\centering
\caption{Campos de la palabra ARINC 429.}
\begin{tabular}{lll}
\toprule
Campo & Bits & Descripción \\
\midrule
Label & 1--8 & Identificador de la variable (octal, transmitido en orden invertido) \\
SDI & 9--10 & \emph{Source/Destination Identifier} \\
Data Field & 11--29 & Dato codificado en BNR o BCD (19 bits) \\
SSM & 30--31 & \emph{Sign/Status Matrix} (estado del dato) \\
Parity & 32 & Paridad impar sobre los 32 bits \\
\bottomrule
\end{tabular}
\end{table}

\subsection{Codificación BNR}
El formato BNR codifica valores numéricos como una fracción del rango máximo:
\begin{equation}
  \text{Valor} = \frac{\text{Data Field}}{2^{18}}\times\text{Escala Máxima}.
\end{equation}
Para codificar un valor positivo:
\begin{equation}
  \text{Data Field} = \text{round}\!\left(\frac{\text{Valor}}{\text{Escala Máxima}}\times(2^{18}-1)\right).
\end{equation}
La resolución resultante es
\begin{equation}
  \text{Resolución} = \frac{\text{Escala Máxima}}{2^{18}}\approx\frac{\text{Escala Máxima}}{262144}.
\end{equation}

\subsection{Matriz SSM (Sign/Status Matrix)}
\begin{table}[H]\centering
\caption{Codificación de la matriz SSM.}
\begin{tabular}{lll}
\toprule
Bits 31--30 & Estado & Significado \\
\midrule
00 & \emph{Failure Warning} & El sensor reporta una falla \\
01 & \emph{No Computed Data} & Dato no disponible o en inicialización \\
10 & \emph{Functional Test} & Dato de prueba, no operacional \\
11 & \emph{Normal Operation} & Dato válido y operacional \\
\bottomrule
\end{tabular}
\end{table}

\subsection{Labels utilizados en esta práctica}
\begin{table}[H]\centering
\caption{Labels ARINC 429 asignados a cada variante.}
\begin{tabular}{lllll}
\toprule
Label (Octal) & Variable & Escala & Unidades & Variante \\
\midrule
206 & \emph{Airspeed} & 0--512 & nudos (kt) & V1 \\
212 & \emph{Temperature} & 0--512 & \si{\celsius} & V2 \\
203 & \emph{Altitude} & 0--131072 & pies (ft) & V3 \\
214 & \emph{Surface Position} & $\pm 180$ & grados ($\si{\degree}$) & V4 \\
\bottomrule
\end{tabular}
\end{table}

\subsection{Simplificaciones de laboratorio}
\begin{table}[H]\centering
\caption{Simplificaciones del estándar ARINC 429 para el entorno de laboratorio.}
\begin{tabular}{lll}
\toprule
Aspecto & ARINC 429 Real & Implementación de Lab \\
\midrule
Medio físico & Par trenzado, $\pm\SI{10}{\volt}$ diferencial & UART 3.3\,V TTL (TX/RX) \\
Velocidad & 12.5 o 100\,kbit/s & 115200 baud (UART) \\
Formato & \emph{Bitstream} serial 32 bits & Cadena hexadecimal ASCII \\
Receptores & Hasta 20 en paralelo & 1 receptor (punto a punto) \\
\bottomrule
\end{tabular}
\end{table}

\subsection{Revisión de conceptos de prácticas anteriores}
\begin{table}[H]\centering
\caption{Relación de conceptos previos con las variantes de P8.}
\begin{tabular}{lll}
\toprule
Práctica & Concepto & Aplicación en P8 \\
\midrule
P1 & GPIO & Indicadores visuales de estado (V2) \\
P2 & ADC/Posición & Sensor de velocidad/altitud/posición (V1, V3, V4) \\
P3 & Temperatura & Sensor de temperatura (V2) \\
P4 & Presión & Fundamento teórico de altimetría (V3) \\
P5 & Servomotores & Indicador de altitud y superficie (V3, V4) \\
P6 & Transistores & Buzzer de alerta (V2), motor (V4) \\
P7 & Motores a pasos & Indicador de aguja mecánica (V1) \\
\bottomrule
\end{tabular}
\end{table}

% ============================================================================
\section{Consideraciones de Seguridad y Buenas Prácticas}
\begin{enumerate}
  \item \textbf{GND común}: ambos MCUs y todas las fuentes de alimentación deben
        compartir la misma referencia de tierra. Sin GND común, la comunicación UART no
        funcionará.
  \item \textbf{Alimentación separada para actuadores}: los servomotores, motores DC y
        motores a pasos deben alimentarse desde fuentes externas (5\,V para servos,
        12\,V para motores DC). No alimentar actuadores desde el pin de 3.3\,V del MCU.
  \item \textbf{Diodos \emph{flyback}}: obligatorios en cargas inductivas (\emph{buzzer},
        motor DC, bobinas del \emph{stepper}). El ULN2003 tiene diodos integrados en el
        pin COM.
  \item \textbf{Protección de entradas ADC}: no aplicar voltajes superiores a 3.3\,V en
        los pines ADC.
  \item \textbf{Fail-safe}: cuando el nodo indicador detecte una palabra con paridad
        inválida o SSM $\ne$ Normal, debe llevar los actuadores a una posición segura.
  \item \textbf{Secuencia de encendido}: encender primero el nodo indicador (que ejecuta
        \emph{homing} si aplica), luego el nodo sensor.
  \item \textbf{Módulo de dron (\S9)}: retirar siempre las hélices para pruebas de banco;
        alimentar los ESC desde una LiPo dedicada, nunca desde el riel del MCU; tener
        siempre accesible la parada de emergencia. Ver \S9.4 para el procedimiento
        completo.
\end{enumerate}

% ============================================================================
\section{Equipo y Materiales}
\subsection{Equipo base (todas las variantes)}
\begin{itemize}
  \item 2$\times$ Microcontroladores (ESP32 o RP2040) — uno por nodo, o 1$\times$ si se
        usa Arquitectura B (lazo interno).
  \item Cable USB para programación.
  \item Protoboard grande y cables \emph{jumper}.
  \item Multímetro digital (DMM).
  \item Osciloscopio (opcional, para verificar señales).
  \item Laptop con IDE (Thonny / VS Code + Pymakr).
\end{itemize}

\subsection{Componentes por variante}
\begin{table}[H]\centering
\caption{Componentes requeridos por cada variante de proyecto.}
\begin{tabular}{lcccc}
\toprule
Componente & V1 & V2 & V3 & V4 \\
\midrule
Potenciómetro \SI{10}{\kilo\ohm} & Sí & --- & Sí & Sí \\
Termistor NTC \SI{10}{\kilo\ohm} & --- & Sí & --- & --- \\
R \SI{10}{\kilo\ohm} (divisor) & --- & Sí & --- & --- \\
Motor 28BYJ-48 + ULN2003 & Sí & --- & --- & --- \\
Servomotor SG90 & --- & --- & Sí & Sí \\
LEDs (verde, amarillo, rojo) + R \SI{330}{\ohm} & --- & Sí & --- & --- \\
Buzzer + 2N2222A + R base \SI{1.2}{\kilo\ohm} & --- & Sí & --- & --- \\
Diodo 1N4007 (\emph{flyback}) & --- & Sí & --- & Sí \\
MOSFET IRF540N + Motor DC & --- & --- & --- & Sí \\
Fin de carrera (\emph{endstop}) & Sí & --- & --- & --- \\
Fuente externa 5--12\,V & Sí & --- & --- & Sí \\
\bottomrule
\end{tabular}
\end{table}
\textbf{Módulo opcional de dron} (\S9): 4$\times$ motor brushless 2212 (900--1400\,KV),
4$\times$ ESC 20--30\,A, 4$\times$ hélice (sin montar en banco), batería LiPo 3S
$\ge$4000\,mAh, cargador balanceador, \emph{frame} clase 250--450, PDB, conectores XT60 y
\emph{bullet} 3.5\,mm — ver \code{docs/dron\_2212.md}.

% ============================================================================
\section{Etapa de Diseño}

\subsection{Análisis de carga eléctrica}
Antes de la integración, se debe calcular la potencia total requerida por el sistema:
\begin{equation}
  P_\text{total} = \sum V_{cc}\cdot I_\text{quiescent} + \sum V_\text{ext}\cdot I_\text{actuador}.
\end{equation}
\begin{table}[H]\centering
\caption{Presupuesto de potencia del sistema integrado.}
\begin{tabular}{lccc}
\toprule
Subsistema & Voltaje (V) & Corriente Pico (mA) & Disipación (mW) \\
\midrule
Nodo Sensor (MCU + ADC) & 3.3 & 50 & 165 \\
Nodo Indicador (MCU) & 3.3 & 50 & 165 \\
Servomotor SG90 & 5.0 & 600 & 3000 \\
Motor a pasos 28BYJ-48 & 5.0 & 250 & 1250 \\
Motor DC (vía MOSFET) & 12.0 & 500 & 6000 \\
Buzzer + LEDs & 3.3 & 80 & 264 \\
\bottomrule
\end{tabular}
\end{table}

\subsection{Codificación BNR — ejemplo de cálculo}
Para la Variante 1 (\emph{Airspeed}), con un valor de 150\,kt y escala máxima de
512\,kt:
\begin{equation}
  \text{Data Field} = \text{round}\!\left(\frac{150}{512}\times(2^{19}-1)\right) =
  \text{round}(153598.8) = 153599.
\end{equation}
En binario: $0010\_0101\_0111\_1111\_111_2$.

\begin{nota}[BNR de 19 bits vs.\ 18 bits con signo]
Para variables exclusivamente positivas (V1, V2, V3), se utilizan los 19 bits completos
del \emph{Data Field} (bits 11--29). Para la Variante 4 (superficie de control con
valores con signo), el bit 29 actúa como bit de signo y se dispone de 18 bits de
magnitud.
\end{nota}

% ============================================================================
\section{Procedimiento Experimental}

\subsection{Arquitectura del sistema}
\subsubsection{Arquitectura A: dos MCUs (distribuida — recomendada)}
\begin{figure}[H]\centering
\begin{tikzpicture}[node distance=3.5cm]
  \node[term] (a) {Nodo Sensor (MCU A)};
  \node[term, right=3.6cm of a] (b) {Nodo Indicador (MCU B)};
  \draw[flow] (a.east) -- node[above,font=\scriptsize]{UART TX $\to$ RX} (b.west);
  \draw[flow, dashed] ([yshift=-4pt]b.west) -- node[below,font=\scriptsize]{GND común} ([yshift=-4pt]a.east);
\end{tikzpicture}
\caption{Arquitectura A: sistema distribuido con dos microcontroladores.}
\end{figure}
Configuración UART: 115200\,baud, 8 bits de datos, sin paridad UART, 1 \emph{stop} bit
(8N1). Conexión cruzada: TX del nodo sensor al RX del nodo indicador, GND común.

\subsubsection{Arquitectura B: un solo MCU (lazo interno)}
Para equipos con un solo MCU disponible, el nodo sensor y el nodo indicador comparten el
mismo microcontrolador. La palabra ARINC 429 se codifica en memoria, se almacena en una
variable de 32 bits y se decodifica inmediatamente. Esta configuración permite validar
la codificación/decodificación sin comunicación serial — es también la arquitectura
recomendada para el \textbf{módulo de análisis de señales} (\S8) y el módulo de
\textbf{dron} (\S9), que corren en un único MCU.

\subsection{Biblioteca ARINC 429 — funciones comunes}
Funciones de referencia (\code{arinc429\_encode}, \code{arinc429\_decode},
\code{arinc429\_send\_uart}, \code{arinc429\_receive\_uart}): calculan el \emph{Data
Field} BNR, empaquetan Label/SDI/SSM, calculan la paridad impar, y transmiten/reciben la
palabra como cadena hexadecimal de 8 caracteres por UART. Ver Apéndice~A.

\subsection{Variante 1: indicador de velocidad aerodinámica (\emph{Airspeed Indicator})}
\textbf{Concepto aeronáutico}: el indicador de velocidad aerodinámica (ASI) es un
instrumento primario de vuelo. El \emph{Air Data Computer} (ADC) codifica la velocidad en
ARINC 429 Label 206 y la envía al \emph{display} de cabina.

\textbf{Adquisición}: potenciómetro de \SI{10}{\kilo\ohm} conectado al ADC, simulando un
transductor de presión dinámica. Mapeo lineal: 0\,V $=$ 0\,kt, 3.3\,V $=$ 300\,kt.
Codificación: Label 206 (octal), escala máxima 512\,kt, BNR positivo, SSM$=11$ (\emph{Normal
Operation}).

\textbf{Salida}: motor a pasos 28BYJ-48 con \emph{driver} ULN2003, posicionando una aguja
proporcional a la velocidad:
\begin{equation}
  N_\text{pasos} = \frac{\text{Velocidad (kt)}}{300}\times\frac{270}{360}\times 4096 =
  \frac{\text{Velocidad (kt)}}{300}\times 3072,
\end{equation}
donde $270\si{\degree}$ es el barrido del dial y 4096 el número de pasos por revolución
en \emph{half-step}.

\begin{table}[H]\centering
\caption{Tabla de verificación — Variante 1 (Airspeed Indicator).}
\begin{tabular}{ccccc}
\toprule
Pot (\%) & Velocidad (kt) & ARINC (hex) & Pasos & Ángulo \\
\midrule
0 & 0 & & 0 & $0\si{\degree}$ \\
25 & 75 & & 768 & $67.5\si{\degree}$ \\
50 & 150 & & 1536 & $135\si{\degree}$ \\
75 & 225 & & 2304 & $202.5\si{\degree}$ \\
100 & 300 & & 3072 & $270\si{\degree}$ \\
\bottomrule
\end{tabular}
\end{table}

\subsection{Variante 2: sistema de alerta de temperatura (\emph{Temperature Alert System})}
\textbf{Concepto aeronáutico}: los sistemas EICAS monitorean continuamente las
temperaturas críticas del motor y generan alertas escalonadas cuando los valores exceden
los límites operativos. El \emph{Air Data Computer} codifica la temperatura en ARINC 429
Label 212.

\textbf{Adquisición}: termistor NTC \SI{10}{\kilo\ohm} ($\beta\approx3950$\,K) en divisor
de voltaje con resistencia fija de \SI{10}{\kilo\ohm}. Conversión mediante la ecuación
Beta: $1/T=1/T_0+(1/\beta)\ln(R_\text{NTC}/R_0)$, con $T_0=298.15$\,K (25\si{\celsius}),
$R_0=\SI{10000}{\ohm}$.

\textbf{Salida}: LED verde ($T<$ umbral amarillo), LED amarillo (advertencia), LED rojo
+ \emph{buzzer} (peligro). Umbrales sugeridos: $T_\text{amarillo}=35\si{\celsius}$,
$T_\text{rojo}=45\si{\celsius}$. El \emph{buzzer} se activa mediante un transistor
2N2222A en configuración \emph{low-side} con diodo \emph{flyback} 1N4007.

\begin{table}[H]\centering
\caption{Tabla de verificación — Variante 2 (Temperature Alert).}
\begin{tabular}{lccccc}
\toprule
Temperatura & SSM & Verde & Amarillo & Rojo & Buzzer \\
\midrule
$25\si{\celsius}$ & Normal & ON & OFF & OFF & OFF \\
$38\si{\celsius}$ & Normal & OFF & ON & OFF & OFF \\
$50\si{\celsius}$ & Normal & OFF & OFF & ON & ON \\
Desconectado & Failure & OFF & ON & ON & OFF \\
\bottomrule
\end{tabular}
\end{table}

\subsection{Variante 3: indicador de altitud barométrica (\emph{Altitude Indicator})}
\textbf{Concepto aeronáutico}: el altímetro barométrico determina la altura de la
aeronave midiendo la presión atmosférica. El \emph{Air Data Computer} calcula la altitud
y la codifica en ARINC 429 Label 203. Un potenciómetro simula la salida del transductor
de presión. Mapeo: 0--3.3\,V $\to$ 0--10000\,ft.

\textbf{Salida}: servomotor SG90 cuya posición es proporcional a la altitud decodificada:
\begin{equation}
  \theta_\text{servo} = \frac{\text{Altitud (ft)}}{\text{Rango máximo (ft)}}\times 180\si{\degree}.
\end{equation}

\begin{table}[H]\centering
\caption{Tabla de verificación — Variante 3 (Altitude Indicator).}
\begin{tabular}{cccc}
\toprule
Pot (\%) & Altitud (ft) & Ángulo Servo & Pulso PWM ($\mu$s) \\
\midrule
0 & 0 & $0\si{\degree}$ & $\sim 500$ \\
25 & 2500 & $45\si{\degree}$ & $\sim 1056$ \\
50 & 5000 & $90\si{\degree}$ & $\sim 1500$ \\
75 & 7500 & $135\si{\degree}$ & $\sim 1944$ \\
100 & 10000 & $180\si{\degree}$ & $\sim 2500$ \\
\bottomrule
\end{tabular}
\end{table}

\subsection{Variante 4: control de superficie de vuelo (Fly-By-Wire)}
\textbf{Concepto aeronáutico}: en sistemas \emph{Fly-By-Wire} (FBW), la palanca de mando
genera señales eléctricas que se codifican digitalmente y se transmiten por el bus de
datos al actuador de la superficie de control. No existe conexión mecánica directa entre
el piloto y la superficie.

\textbf{Adquisición}: potenciómetro centrado. ADC centrado: 0\,V$=-30\si{\degree}$,
1.65\,V$=0\si{\degree}$ (\emph{neutral}), 3.3\,V$=+30\si{\degree}$. Para valores con
signo, el bit 29 del \emph{Data Field} actúa como signo:
\begin{equation}
  \text{Magnitud BNR} = \text{round}\!\left(\frac{|\text{Deflexión}|}{180}\times(2^{18}-1)\right).
\end{equation}

\textbf{Salida}: servomotor SG90 representando la superficie de elevador. \emph{Neutral}
$=90\si{\degree}$, deflexión $+30\si{\degree}=120\si{\degree}$, deflexión
$-30\si{\degree}=60\si{\degree}$. Actuador secundario (opcional): motor DC controlado por
MOSFET, representando empuje proporcional a deflexión positiva.

\begin{table}[H]\centering
\caption{Tabla de verificación — Variante 4 (Fly-By-Wire).}
\begin{tabular}{ccccc}
\toprule
Pot (\%) & Deflexión & Servo & Motor \\
\midrule
0 & $-30\si{\degree}$ & $60\si{\degree}$ & OFF \\
25 & $-15\si{\degree}$ & $75\si{\degree}$ & OFF \\
50 & $0\si{\degree}$ (\emph{neutral}) & $90\si{\degree}$ & OFF \\
75 & $+15\si{\degree}$ & $105\si{\degree}$ & \(\sim\)50\,\% \\
100 & $+30\si{\degree}$ & $120\si{\degree}$ & 100\,\% \\
\bottomrule
\end{tabular}
\end{table}

\subsection{Caso de estudio: inyección de fallos y validación de integridad}
Independientemente de la variante seleccionada, cada equipo debe realizar las siguientes
pruebas de robustez:
\begin{table}[H]\centering
\caption{Matriz de pruebas de inyección de fallos.}
\begin{tabular}{llll}
\toprule
Condición de Fallo & Procedimiento & Resultado Esperado & Pass \\
\midrule
Paridad corrupta & Invertir bit 32 en palabra & Rechazo, \emph{fail-safe} & \\
SSM = Failure (00) & Forzar ssm=0x00 en sensor & Actuador a posición segura & \\
SSM = NCD (01) & Forzar ssm=0x01 en sensor & Actuador congelado o seguro & \\
Label incorrecto & Enviar Label 0o377 & Ignorar trama & \\
Desconexión UART & Retirar cable TX & \emph{Fail-safe} tras \emph{timeout} & \\
\bottomrule
\end{tabular}
\end{table}

\begin{nota}[Contexto aeronáutico: DO-178C]
En el desarrollo de software aeronáutico, la norma DO-178C (\emph{Software
Considerations in Airborne Systems and Equipment Certification}) establece que los
sistemas de nivel DAL A deben demostrar cobertura completa de código y robustez ante
condiciones anormales. La inyección de fallos es una técnica estándar de verificación.
\end{nota}

% ============================================================================
\section{Análisis de Señales}
\label{sec:senales}

Además de la codificación ARINC 429, la implementación de referencia del repositorio
añade un modo dedicado (\textbf{Modo 9}) para analizar la cadena de señal completa como
un sistema de adquisición--procesamiento--actuación, y aplica los mismos paradigmas de
las Prácticas 4--7 (muestreo, jitter, respuesta dinámica) al sistema integrado.

\subsection{Latencia de la cadena sensor $\to$ actuador}
El Modo 9$\to$1 mide, con \code{ticks\_us}, el tiempo de \emph{leer ADC (altitud)
$\to$ calcular $\to$ mover servo (alerón)} en 200 ciclos y reporta media, $\sigma$ y la
tasa de actualización efectiva. Compare esta latencia con los requisitos típicos de un
sistema FBW primario ($<\SI{50}{\milli\second}$, \S10.2 del registro de datos original) y
discuta si el sistema es apto para control primario de vuelo o solo para monitoreo de
sistemas secundarios.

\subsection{Muestreo multicanal y anti-aliasing}
El Modo 9$\to$2 muestrea los 4 sensores ADC (altitud, velocidad, actitud, luz) a una
$F_s$ común y emite un CSV multicanal. Preguntas a resolver: ¿los 4 canales se leen
realmente al mismo instante o hay un \emph{skew} por leerlos secuencialmente? ¿Qué
$F_s$ es adecuada dado el ancho de banda físico de cada variable simulada? Con el
generador de funciones inyectando un seno conocido en un canal (como en P5, Modo 6), se
puede verificar el aliasing también aquí.

\subsection{Jitter del PWM de los ESC (módulo de dron)}
Si el módulo opcional de dron (\S\ref{sec:dron}, más abajo) está habilitado, el Modo
8$\to$4 mide el \emph{jitter} de la señal PWM que llega a un ESC, con el mismo método que
el Modo 5 de la Práctica 5 (osciloscopio, \emph{Pulse Width $\to$ StdDev}).

% ============================================================================
\section{Módulo Opcional: Dron con Motores 2212}
\label{sec:dron}

La implementación de referencia incluye un \textbf{Modo 8}, deshabilitado por defecto
(\code{ENABLE\_DRONE = False}), que añade la capa de actuación de bajo nivel para
construir un cuadricóptero real con 4 motores brushless \textbf{2212} (900--1400\,KV) y
sus ESC. Este módulo se documenta en detalle en \code{docs/dron\_2212.md}; aquí se resume
lo esencial.

\subsection{Driver de ESC con armado fail-safe}
Un ESC recibe la misma señal que un servo (PWM 50\,Hz, 1000--2000\,\si{\micro\second}),
pero la interpreta como \emph{throttle} (0--100\,\%) y exige una secuencia de
\textbf{armado}: throttle mínimo sostenido $\sim$2\,s antes de aceptar comandos. El
\emph{driver} del repositorio (\code{lib/esc.py} / \code{common/esc.h}) implementa una
máquina de estados que \textbf{niega el throttle} si no se armó antes — un fallo seguro
por diseño.

\subsection{Mezclador de cuadricóptero en X}
El módulo \code{lib/quad\_mixer.py} / \code{common/quad\_mixer.h} traduce comandos de
throttle/roll/pitch/yaw a 4 valores de throttle por motor:
\begin{align}
  M_1 &= \text{throttle}+\text{roll}-\text{pitch}+\text{yaw} \quad (\text{front-left, CW}), \\
  M_2 &= \text{throttle}-\text{roll}-\text{pitch}-\text{yaw} \quad (\text{front-right, CCW}), \\
  M_3 &= \text{throttle}+\text{roll}+\text{pitch}-\text{yaw} \quad (\text{rear-left, CCW}), \\
  M_4 &= \text{throttle}-\text{roll}+\text{pitch}+\text{yaw} \quad (\text{rear-right, CW}).
\end{align}

\begin{nota}[Lazo abierto — ruta de extensión]
Este mezclador es de \textbf{lazo abierto}: no hay IMU ni PID, así que \textbf{no
estabiliza un vuelo real} por sí solo. Es la capa de actuación (throttle $\to$ señal ESC)
sobre la que se puede construir un controlador de vuelo completo añadiendo un IMU
(p.\ ej.\ MPU6050), un filtro de fusión y un lazo PID por eje — ver
\code{docs/dron\_2212.md}, \S7, para la ruta de extensión detallada. Esto queda fuera del
alcance evaluado de esta práctica.
\end{nota}

\subsection{Submenú del Modo 8}
\begin{enumerate}
  \item \textbf{Armar los 4 ESC} — requiere confirmación explícita de que las hélices
        están retiradas.
  \item \textbf{Test individual de motor} — throttle manual bajo (\(\le\)20\,\%
        recomendado sin hélice).
  \item \textbf{Mezclador X} — throttle/roll/pitch/yaw por teclado, con parada de
        emergencia en la barra espaciadora.
  \item \textbf{Jitter del PWM del ESC} — osciloscopio, igual método que P5 Modo 5.
  \item \textbf{Parada de emergencia} — throttle 0 en los 4 motores y desarmado
        inmediato.
\end{enumerate}

\subsection{Seguridad (resumen — ver docs/dron\_2212.md)}
Retirar siempre las hélices para pruebas de banco. Alimentar los ESC desde una LiPo 3S
dedicada (nunca el riel del MCU). Mantener la parada de emergencia siempre accesible.
Verificar el sentido de giro de cada motor antes de considerar el montaje de hélices
(fuera del alcance de laboratorio de esta práctica).

% ============================================================================
\section{Alternativa de Trabajo en Casa (Multímetro + Simulación)}
\label{sec:casa}

\paragraph{Con multímetro (en casa).} Verificar el voltaje de la LiPo (11.1--12.6\,V
para 3S) antes de conectar el módulo de dron. Verificar continuidad de los conectores
XT60/\emph{bullet} y que el GND es común entre MCU y ESC. Para las 4 variantes
ARINC 429, medir el voltaje del sensor (potenciómetro/NTC) y calcular a mano el
\emph{Data Field} esperado, comparándolo con el valor hexadecimal que imprime el
firmware.

\paragraph{Con simulación.} El Modo 9 (latencia y muestreo multicanal) y toda la
codificación/decodificación BNR (Arquitectura B, un solo MCU) son \textbf{100\,\%
software}: no requieren osciloscopio, generador ni un segundo MCU, y pueden ejecutarse
íntegramente en casa. Para quien no tenga 2 MCU disponibles, la Arquitectura B es la
recomendada. La lógica de codificación/decodificación (Apéndice~A) también puede
replicarse en un intérprete de Python en la PC para verificarla sin ningún hardware.

\paragraph{Requiere laboratorio presencial.} El hardware de cada variante (motor a
pasos, servo, \emph{buzzer}, MOSFET+motor DC) y el módulo opcional de dron requieren los
componentes físicos; el jitter del PWM del ESC (Modo 8$\to$4) requiere además
osciloscopio.

\begin{table}[H]\centering
\caption{Alternativas de trabajo en casa.}
\begin{tabular}{lccc}
\toprule
Actividad & Multímetro & Simulación & Requiere laboratorio \\
\midrule
Codificación/decodificación BNR & --- & Sí (100\,\%) & No \\
Variantes 1--4 (hardware real) & Sí (sensor) & Parcial & Sí \\
Inyección de fallos (paridad/SSM) & --- & Sí (lógica) & No \\
Modo 9 (latencia, muestreo) & --- & --- & No (usa el propio MCU) \\
Modo 8 (dron, armar/mezclador) & Sí (LiPo, continuidad) & --- & Sí (motores 2212) \\
Modo 8 (jitter ESC) & No & No (idealizada) & \textbf{Sí, osciloscopio} \\
\bottomrule
\end{tabular}
\end{table}

% ============================================================================
\section{Registro de Datos y Comparativa}

\subsection{Tabla de validación integral}
\begin{table}[H]\centering
\caption{Registro de validación integral — completar durante el laboratorio.}
\begin{tabular}{lccccc}
\toprule
Entrada & Calculado & ARINC (hex) & Decodificado & Error (\%) & SSM \\
\midrule
 & & & & & \\
 & & & & & \\
 & & & & & \\
\bottomrule
\end{tabular}
\end{table}

\subsection{Análisis de latencia}
Registre el tiempo de respuesta del sistema desde el cambio en el sensor hasta la
respuesta visible del actuador (o use directamente el Modo 9$\to$1). Compare con los
requisitos típicos de un sistema FBW primario ($<\SI{50}{\milli\second}$).

% ============================================================================
\section{Análisis de Resultados}
\begin{enumerate}
  \item \textbf{Validación de codificación}: calcule el porcentaje de error entre el
        valor físico original y el valor decodificado. Relacione el error con la
        resolución BNR de la ecuación 3.
  \item \textbf{Efectividad de la detección de fallos}: discuta cómo la paridad y el
        SSM evitaron movimientos no comandados durante las pruebas de inyección de
        fallos.
  \item \textbf{Análisis de latencia}: grafique el tiempo de respuesta (Modo 9$\to$1) y
        determine si es adecuado para control primario de vuelo o solo para monitoreo de
        sistemas secundarios.
  \item \textbf{Comparativa con ARINC 429 real}: calcule la tasa de transmisión efectiva
        de su sistema en palabras/segundo y compárela con la tasa típica de ARINC 429 a
        alta velocidad ($\sim$381\,palabras/s).
  \item \textbf{Tolerancia de componentes}: evalúe cómo la tolerancia del potenciómetro
        (5\,\%) y la no linealidad del ADC afectan el valor BNR codificado.
  \item \textbf{Muestreo multicanal} (Modo 9$\to$2): ¿los 4 canales presentan
        \emph{skew} detectable? ¿La $F_s$ elegida es suficiente para la dinámica de
        cada variable simulada?
\end{enumerate}

% ============================================================================
\section{Preguntas de Discusión y Refuerzo}
\begin{enumerate}
  \item Por qué ARINC 429 utiliza paridad impar en lugar de paridad par? ¿Qué ventaja
        ofrece para detectar errores de un solo bit?
  \item Si el bus ARINC 429 transporta una palabra de altitud con paridad correcta pero
        el SSM indica ``Failure Warning'', ¿debe el indicador mostrar el valor o
        ignorarlo? Justifique considerando factores humanos.
  \item ¿Cuál es la resolución angular lograda por el motor a pasos 28BYJ-48 en
        \emph{half-step} para la Variante 1? ¿Es suficiente para un indicador de
        velocidad aerodinámica real?
  \item En la Variante 4 (Fly-By-Wire), ¿qué ocurriría si el nodo actuador pierde
        comunicación con el nodo sensor durante 500\,ms? ¿Cuál sería la estrategia de
        \emph{fail-safe} según DO-178C?
  \item Compare la latencia de su sistema (Modo 9$\to$1) con los requisitos típicos de
        un sistema FBW primario ($<\SI{50}{\milli\second}$). ¿Qué factores contribuyen a
        la latencia?
  \item ¿Por qué se utiliza una cadena hexadecimal ASCII sobre UART en lugar de
        transmitir los 4 \emph{bytes} directamente?
  \item Explique cómo la arquitectura unidireccional de ARINC 429 simplifica el diseño
        del bus frente a protocolos bidireccionales como CAN o ARINC 664 (AFDX).
  \item Si se quisiera añadir redundancia dual al sistema, ¿cómo modificaría la
        arquitectura? Describa el concepto de \emph{voting} en aviónica.
  \item ¿Qué mejoras implementaría para que su sistema cumpla con DAL A según DO-178C?
  \item Calcule la tasa de transmisión efectiva de su sistema en palabras/segundo y
        compárela con ARINC 429 a alta velocidad.
  \item \textbf{(Módulo de dron)} ¿Por qué el mezclador X es de lazo abierto y qué
        componentes (sensor + algoritmo) le faltan para estabilizar un vuelo real? ¿Qué
        norma o buena práctica de la industria de drones establece los requisitos
        mínimos de \emph{failsafe} de radio para un vehículo no tripulado?
\end{enumerate}

% ============================================================================
\section{Entregables (Proyecto Final — Formato AIAA)}
Cada equipo debe entregar:
\begin{enumerate}
  \item Reporte técnico (formato AIAA) que incluya: resumen ejecutivo, diagrama de
        arquitectura del sistema, esquemático eléctrico completo, tabla de asignación
        de pines, tabla de codificación ARINC 429, memoria de cálculo, resultados con
        porcentaje de error, análisis de latencia, prueba de inyección de fallos, y
        respuestas a las preguntas de discusión.
  \item Código fuente completo y comentado (ambos nodos, o el panel integrado).
  \item Video o evidencia fotográfica del sistema en operación, mostrando: \emph{homing}
        del motor a pasos (si aplica), respuesta del actuador ante variaciones del
        sensor, y respuesta ante inyección de fallo.
  \item \textbf{(Si se habilitó el módulo de dron)} evidencia del armado seguro de los 4
        ESC y del mezclador X operando \textbf{sin hélices}.
\end{enumerate}

% ============================================================================
\section{Referencias}
\begin{enumerate}
  \item ARINC Specification 429 Part 1-18: Mark 33 Digital Information Transfer System.
        Aeronautical Radio, Inc.
  \item RTCA DO-178C: Software Considerations in Airborne Systems and Equipment
        Certification. 2011.
  \item RTCA DO-254: Design Assurance Guidance for Airborne Electronic Hardware. 2000.
  \item Moir, I., \& Seabridge, A. (2013). \emph{Aircraft Systems: Mechanical,
        Electrical, and Avionics Subsystems Integration}. 4th Ed., Wiley.
  \item Spitzer, C.R. (2015). \emph{Digital Avionics Handbook}. 3rd Ed., CRC Press.
  \item Collinson, R.P.G. (2011). \emph{Introduction to Avionics Systems}. 3rd Ed.,
        Springer.
  \item \emph{Datasheets} de motores 2212 (fabricante según modelo adquirido); BLHeli /
        BLHeli\_S ESC Documentation — ver \code{docs/dron\_2212.md} para la lista
        completa del módulo opcional.
\end{enumerate}

% ============================================================================
\appendix
\section{Código Fuente del Repositorio}
\begin{table}[H]\centering
\caption{Estructura de archivos de la Práctica 8.}
\begin{tabular}{ll}
\toprule
Ruta & Descripción \\
\midrule
\code{MicroPython/\{ESP32,RP2040\}/P8/main.py} & Panel integrado (9 modos) \\
\code{.../P8/lib/\{sensors,flight\_controls,propulsion,landing\_gear\}.py} & Subsistemas base \\
\code{.../P8/lib/\{esc,quad\_mixer\}.py} & Módulo opcional de dron \\
\code{.../P8/lib/siglab.py} & Utilidades de análisis de señales (Modo 9) \\
\code{.../P8/docs/dron\_2212.md} & Guía completa del módulo opcional \\
\code{.../P8/docs/analisis\_senales.md} & Guía del Modo 9 \\
\code{C++/SISELA-CPP/src/practices/p8.cpp} & Implementación C++ (9 modos) \\
\code{C++/.../common/\{esc,quad\_mixer,siglab\}.h} & Módulos comunes C++ \\
\code{tools/sisela\_signal/} & Toolkit PC de análisis de señales \\
\bottomrule
\end{tabular}
\end{table}

\section{Contexto Aeronáutico Complementario}
\textbf{Sistemas reales de aviónica integrada}: el FMS (\emph{Flight Management System})
integra datos de múltiples sensores (ADC, IRS, GPS) vía ARINC 429/629 para calcular la
trayectoria óptima. El EICAS (\emph{Engine Indication and Crew Alerting System}) recibe
docenas de palabras ARINC 429 de los motores y genera alertas jerarquizadas (\emph{Warning},
\emph{Caution}, \emph{Advisory}) — el mismo patrón que la Variante 2 de esta práctica. El
EFIS (\emph{Electronic Flight Instrument System}) reemplaza los instrumentos analógicos
mecánicos por pantallas que decodifican ARINC 429 en tiempo real, aunque muchos aviones
de aviación general conservan instrumentos de respaldo mecánicos (motor a pasos o
servomotor) idénticos en principio a los de esta práctica.

\end{document}
